%! Simplifying Radicals; Adding \& Subtracting — Unit 6, Lesson 6.2 — Slides (Beamer)
\documentclass[aspectratio=169, 11pt]{beamer}
\usepackage{algebra2-beamer}   % provides \forestheader{} and \sectionlabel[color]{}

\begin{document}

% TITLE SLIDE (CourseName is not defined in beamer, so the name is literal)
{
\setbeamercolor{background canvas}{bg=forest}
\begin{frame}[plain]
  \vspace{2.8cm}\hspace{0.55cm}%
  \begin{minipage}{0.9\textwidth}
    {\color{goldacc}\bfseries\small LESSON 6.2}\\[0.5cm]
    {\color{white}\bfseries\LARGE Simplifying Radicals; Adding \& Subtracting}\\[1.2cm]
    {\color{forestmid}\small Algebra 2: Shepherd~$\cdot$~Unit 6: Radical Functions}
  \end{minipage}
\end{frame}
}

% HOOK
\begin{frame}[t, plain]
  \forestheader{Two of these are true. One is a lie.}
  \sectionlabel{Hook}
  \begin{columns}[T]
    \begin{column}{0.52\textwidth}
      \begin{itemize}\setlength{\itemsep}{8pt}
        \item[(a)] $\sqrt{8}\cdot\sqrt{2} = 4$
        \item[(b)] $\sqrt{9}+\sqrt{16} = \sqrt{25}$
        \item[(c)] $\sqrt{8}+\sqrt{18} = 5\sqrt{2}$
      \end{itemize}
      \vspace{6pt}
      {\footnotesize Test each one with decimals. Which one breaks?}
    \end{column}
    \begin{column}{0.44\textwidth}
      \begin{block}{The verdict}
        (a) \textbf{true} --- $\sqrt{16}=4$ \\[3pt]
        (b) \textcolor{keyred}{\textbf{false}} --- $3+4=7$, not $5$ \\[3pt]
        (c) \textbf{true} --- $2.828+4.243=7.071$
      \end{block}
      \vspace{3pt}
      {\footnotesize The true ones involve \textbf{multiplication} inside the radical. The lie tries
      to split over \textbf{addition}.}
    \end{column}
  \end{columns}
\end{frame}

% WHERE THE PROPERTIES COME FROM
\begin{frame}[t, plain]
  \forestheader{The properties are yesterday's rules in a costume}
  \sectionlabel[goldacc]{Product \& Quotient}
  \begin{columns}[T]
    \begin{column}{0.48\textwidth}
      \[ \sqrt[n]{ab} = (ab)^{1/n} = a^{1/n}b^{1/n} = \sqrt[n]{a}\cdot\sqrt[n]{b} \]
      \vspace{-6pt}
      \begin{block}{The two properties}
        \centering
        $\displaystyle \sqrt[n]{ab}=\sqrt[n]{a}\cdot\sqrt[n]{b}$ \\[5pt]
        $\displaystyle \sqrt[n]{\frac{a}{b}}=\frac{\sqrt[n]{a}}{\sqrt[n]{b}}$ \quad $(b\ne0)$
      \end{block}
    \end{column}
    \begin{column}{0.48\textwidth}
      \footnotesize
      \textbf{The fine print.} For an \textbf{even} index, $a$ and $b$ must be non-negative --- or
      the pieces are not real numbers. You may \emph{not} write
      $\sqrt{-8}=\sqrt{-4}\cdot\sqrt{2}$.
      \vspace{6pt}
      \begin{block}{And there is no third property}
        $\sqrt{9\cdot16}=12$ \; and \; $\sqrt9\cdot\sqrt{16}=12$ \quad \textcolor{forest}{\textbf{$\checkmark$}} \\[3pt]
        $\sqrt{9+16}=5$ \; but \; $\sqrt9+\sqrt{16}=7$ \quad \textcolor{keyred}{$\times$}
      \end{block}
      \vspace{3pt}
      {\footnotesize \textcolor{goldacc}{\textbf{$\sqrt{a+b}\ne\sqrt a+\sqrt b$. Ever.}}}
    \end{column}
  \end{columns}
\end{frame}

% SIMPLEST FORM
\begin{frame}[t, plain]
  \forestheader{Simplest radical form --- and pulling out the \emph{largest} power}
  \sectionlabel{Simplifying}
  \begin{columns}[T]
    \begin{column}{0.44\textwidth}
      \begin{block}{A radical is simplified when}
        \footnotesize
        1. no factor of the radicand is a perfect $n$th power; \\[3pt]
        2. there is no fraction under the radical; \\[3pt]
        3. there is no radical in a denominator. \\[3pt]
        {\scriptsize (Condition 3 is Lesson 6.3.)}
      \end{block}
    \end{column}
    \begin{column}{0.52\textwidth}
      \[ \sqrt{72}=\sqrt{36\cdot2}=\sqrt{36}\cdot\sqrt2=6\sqrt2 \]
      \vspace{-8pt}
      \footnotesize
      \begin{itemize}\setlength{\itemsep}{4pt}
        \item Spot only $72=4\cdot18$? You get $2\sqrt{18}$ --- \textbf{true but not finished}.
        \item $2\sqrt{18}=2\cdot3\sqrt2=6\sqrt2$. Same answer, more steps.
        \item \textcolor{goldacc}{\textbf{Always ask last: is there still a perfect square in
              there?}}
      \end{itemize}
      \vspace{4pt}
      {\footnotesize Backup: prime-factor it. $72=2\cdot2\cdot2\cdot3\cdot3$ --- one factor escapes
      per \emph{pair}.}
    \end{column}
  \end{columns}
\end{frame}

% HIGHER INDICES
\begin{frame}[t, plain]
  \forestheader{The index sets the group size}
  \sectionlabel[goldacc]{Higher indices}
  \begin{columns}[T]
    \begin{column}{0.50\textwidth}
      \footnotesize
      \begin{tabular}{ll}
      \hline
      Squares & $4,9,16,25,36,49,64,81,100,121,144$ \\
      Cubes & $8,27,64,125,216$ \; (and $-8,-27,\dots$) \\
      4th powers & $16,81,256$ \qquad 5th: $32,243$ \\
      \hline
      \end{tabular}
      \vspace{8pt}

      \begin{itemize}\setlength{\itemsep}{4pt}
        \item $\sqrt[3]{54}=\sqrt[3]{27\cdot2}=3\sqrt[3]{2}$
        \item $\sqrt[4]{48}=\sqrt[4]{16\cdot3}=2\sqrt[4]{3}$
        \item $\sqrt[3]{-24}=-2\sqrt[3]{3}$ \; {\scriptsize (odd index, negative is fine)}
      \end{itemize}
    \end{column}
    \begin{column}{0.46\textwidth}
      \begin{block}{The trap}
        Under a cube root, a \emph{pair} is worth nothing. \\[4pt]
        $\sqrt[3]{9}$ is \textbf{already simplified} --- even though $9$ is a perfect square. \\[4pt]
        You need \textbf{three} equal factors to get one out.
      \end{block}
      \vspace{3pt}
      {\footnotesize Squares escape a square root in \textbf{pairs}; cubes escape a cube root in
      \textbf{threes}. The index tells you which.}
    \end{column}
  \end{columns}
\end{frame}

% ALGEBRAIC RADICANDS
\begin{frame}[t, plain]
  \forestheader{Variables: divide the exponent by the index}
  \sectionlabel{Algebraic radicands}
  \begin{columns}[T]
    \begin{column}{0.46\textwidth}
      \[ \sqrt{x^{7}}=\sqrt{x^{6}\cdot x}=x^{3}\sqrt{x} \]
      \vspace{-6pt}
      \begin{block}{The rule}
        \centering Divide the exponent by the index. \\[3pt]
        \textbf{Quotient} goes outside. \\
        \textbf{Remainder} stays inside.
      \end{block}
      \vspace{3pt}
      \footnotesize
      $\sqrt{y^{9}}=y^{4}\sqrt y$ \qquad $\sqrt[3]{x^{11}}=x^{3}\sqrt[3]{x^{2}}$
    \end{column}
    \begin{column}{0.50\textwidth}
      \begin{block}{Why --- and it should look familiar}
        \centering
        $\displaystyle \sqrt{x^{7}}=x^{7/2}=x^{3\frac12}=x^{3}\cdot x^{1/2}=x^{3}\sqrt x$
      \end{block}
      \vspace{4pt}
      {\footnotesize \textcolor{goldacc}{\textbf{Simplifying a radical is rewriting an improper
      fraction as a mixed number.}} The whole-number part comes out front; the leftover fraction
      stays under the radical.}
      \vspace{5pt}

      {\footnotesize All together: \; $\sqrt{50x^{5}y^{8}}=5x^{2}y^{4}\sqrt{2x}$ \\[3pt]
      \qquad\qquad\quad $\sqrt[3]{40a^{7}b^{3}}=2a^{2}b\sqrt[3]{5a}$}
    \end{column}
  \end{columns}
\end{frame}

% QUOTIENT PROPERTY
\begin{frame}[t, plain]
  \forestheader{The quotient property --- and one thing we cannot finish yet}
  \sectionlabel{Fractions}
  \begin{columns}[T]
    \begin{column}{0.50\textwidth}
      \footnotesize
      \begin{itemize}\setlength{\itemsep}{6pt}
        \item $\sqrt{\dfrac{25}{36}}=\dfrac{5}{6}$ \qquad
              $\sqrt[3]{\dfrac{8}{27}}=\dfrac{2}{3}$
        \item $\sqrt{\dfrac{7}{16}}=\dfrac{\sqrt7}{4}$
        \item Backwards: $\dfrac{\sqrt{50}}{\sqrt2}=\sqrt{25}=5$ \\[3pt]
              \qquad\qquad\quad $\dfrac{\sqrt{27}}{\sqrt3}=\sqrt9=3$
      \end{itemize}
    \end{column}
    \begin{column}{0.46\textwidth}
      \begin{block}{Unfinished business}
        $\sqrt{\dfrac{3}{5}}=\dfrac{\sqrt3}{\sqrt5}$ \\[4pt]
        {\footnotesize A radical in the \textbf{denominator} breaks simplest-form condition 3.}
      \end{block}
      \vspace{3pt}
      {\footnotesize Clearing it is called \textbf{rationalizing} --- the first thing we do in
      Lesson 6.3.}
    \end{column}
  \end{columns}
\end{frame}

% LIKE RADICALS
\begin{frame}[t, plain]
  \forestheader{Adding radicals is collecting like terms}
  \sectionlabel[goldacc]{Like radicals}
  \begin{columns}[T]
    \begin{column}{0.48\textwidth}
      \footnotesize
      \begin{tabular}{ll}
      \hline
      $3x+7x=10x$ & $3\sqrt5+7\sqrt5=10\sqrt5$ \\
      $3x+7y$ \; no & $3\sqrt5+7\sqrt2$ \; no \\
      \hline
      \end{tabular}
      \vspace{8pt}

      \begin{block}{Like radicals}
        Same \textbf{index} \emph{and} same \textbf{radicand}.
      \end{block}
      \vspace{4pt}
      {\footnotesize Only the numbers in front change. $3\sqrt5+7\sqrt5$ is $10\sqrt5$ ---
      \textcolor{keyred}{\textbf{not $10\sqrt{10}$}}.}
    \end{column}
    \begin{column}{0.48\textwidth}
      \begin{block}{The catch}
        \centering $\sqrt{8}+\sqrt{18} = 2\sqrt2+3\sqrt2 = 5\sqrt2$
      \end{block}
      \vspace{4pt}
      {\footnotesize Neither term \emph{looked} like the other.
      \textcolor{goldacc}{\textbf{Simplify first; decide second.}} Never say ``cannot combine''
      until both radicals are in simplest form.}
      \vspace{5pt}

      \footnotesize
      $\sqrt{45}-\sqrt{20}=3\sqrt5-2\sqrt5=\sqrt5$ \\[3pt]
      $5\sqrt{12}+2\sqrt{75}-\sqrt3=19\sqrt3$ \\[3pt]
      $\sqrt[3]{2}+\sqrt{2}$ --- the \emph{indices} differ, so no.
    \end{column}
  \end{columns}
\end{frame}

% ERROR SLIDE
\begin{frame}[t, plain]
  \forestheader{The four ways this goes wrong}
  \sectionlabel[goldacc]{Watch out}
  \begin{columns}[T]
    \begin{column}{0.48\textwidth}
      \footnotesize
      \begin{block}{1. Stopping early}
        $\sqrt{80}=2\sqrt{20}$ --- true, but $20$ still has a $4$ in it. Answer: $4\sqrt5$.
      \end{block}
      \vspace{4pt}
      \begin{block}{2. Splitting over addition}
        $\sqrt{20}+\sqrt5\ne\sqrt{25}$. \; Check it: $4.47+2.24=6.71$, not $5$.
      \end{block}
    \end{column}
    \begin{column}{0.48\textwidth}
      \footnotesize
      \begin{block}{3. Pairs under a cube root}
        $\sqrt[3]{54}=\sqrt[3]{9\cdot6}$ gives a $3$ that has no right to leave. \;
        $54=27\cdot2$, so $3\sqrt[3]{2}$.
      \end{block}
      \vspace{4pt}
      \begin{block}{4. Adding the radicands too}
        $3\sqrt7+2\sqrt7=5\sqrt7$, \textbf{not} $5\sqrt{14}$. The radical never changes.
      \end{block}
    \end{column}
  \end{columns}
\end{frame}

% PROOF SLIDE
\begin{frame}[t, plain]
  \forestheader{Why $\sqrt{a}+\sqrt{b}=\sqrt{a+b}$ is \emph{never} an accident}
  \sectionlabel{Extension}
  \begin{columns}[T]
    \begin{column}{0.52\textwidth}
      Suppose it were true, with $a,b\ge0$:
      \[ \sqrt a+\sqrt b=\sqrt{a+b} \]
      Square both sides:
      \[ a+2\sqrt{ab}+b = a+b \]
      \[ \Rightarrow\quad 2\sqrt{ab}=0 \quad\Rightarrow\quad ab=0 \]
    \end{column}
    \begin{column}{0.44\textwidth}
      \begin{block}{Conclusion}
        The two sides agree \textbf{only} when $a=0$ or $b=0$. For any two \emph{positive} numbers
        it is false --- always.
      \end{block}
      \vspace{4pt}
      {\footnotesize The term the mistake throws away is exactly $2\sqrt{ab}$. \\[4pt]
      Check: $a=9$, $b=16$ gives $7$ vs.\ $5$, and the missing piece is
      $2\sqrt{144}=24$ --- under the radical, that is the whole gap.}
    \end{column}
  \end{columns}
\end{frame}

% ROADMAP
\begin{frame}[t, plain]
  \forestheader{Where Unit 6 goes next}
  \sectionlabel{Connections}
  \textbf{Today:} $\sqrt[n]{ab}=\sqrt[n]a\cdot\sqrt[n]b$ pulls perfect powers \emph{out}; like
  radicals then add exactly the way like terms do --- once you have simplified far enough to see
  which ones are alike.\\[8pt]
  \begin{itemize}\setlength{\itemsep}{3pt}
    \item \textbf{6.3} The same property run \emph{forwards}: multiplying, distributing, FOIL, and
          rationalizing a denominator with a \textbf{conjugate}.
    \item \textbf{6.4} Graphing $y=a\sqrt{x-h}+k$ --- where simplifying $\sqrt{8x}$ into
          $2\sqrt{2x}$ turns out to be a \emph{vertical stretch}.
    \item \textbf{6.5} Solving radical equations, and why extraneous solutions appear.
    \item \textbf{6.6--6.7} Composition and inverses --- the unit capstone.
  \end{itemize}
\end{frame}

\end{document}
